\section{The sixfold theorem and its consequences}\label{sec:conclusion}

\begin{proof}[Proof of Theorem~\ref{thm:main}]
Let $X$ satisfy its hypotheses. Since $X$ is projective and positive
dimensional, an ample divisor has nonzero numerical class, so
$\rho(X)\geq1$. If $\rho(X)\geq2$, then
$\rho(X)>\dim X-5$. Kanemitsu's Theorem~0.2
\cite{Kanemitsu2016} applies and gives rational homogeneity.

Suppose now that $\rho(X)=1$.
By Proposition~\ref{prop:pseudoindex-reduction}, either $X$ is
$\mathbb P^6$ or $Q^6$, or its pseudoindex is four or five.
The first two varieties are rational homogeneous.
Proposition~\ref{prop:pseudoindex-four} settles pseudoindex four.
In pseudoindex five, the characteristic-class identities of
Section~\ref{sec:characteristic-classes} and
Lemma~\ref{lem:finite-ratio} give $32-20u+3u^2=0$.
Proposition~\ref{prop:pseudoindex-five-recognition} therefore gives
$X\simeq\operatorname{Gr}(2,5)$, which is rational homogeneous.
These cases exhaust the possibilities.
\end{proof}

\subsection{A K\"ahler--Einstein consequence}

\begin{corollary}\label{cor:positive-ke}
Let $X$ be a smooth connected complex projective Fano variety with nef
tangent bundle and $1\leq\dim_{\mathbb C}X\leq6$.
Then $X$ admits a positive K\"ahler--Einstein metric. More precisely,
with the Ricci-form convention
\[
 \rho(\omega)=-i\partial\bar\partial\log\det(g_{a\bar b}),
\]
there is a K\"ahler form $\omega$ such that
\[
 \rho(\omega)=\omega,\qquad [\omega]=2\pi c_1(-K_X).
\]
\end{corollary}

\begin{proof}
In dimension six, Theorem~\ref{thm:main} makes $X$ rational homogeneous.
For $1\leq n=\dim_{\mathbb C}X\leq5$, its Picard number is at least
one and hence greater than $n-5$, so Kanemitsu's Theorem~0.2
\cite{Kanemitsu2016} gives the same conclusion. Thus the
lower-dimensional cases already follow from the known classification.
The homogeneous metric construction of Alekseevskii and Perelomov
\cite{AlekseevskiiPerelomov1986}, recalled in
\cite[\S1, p.~285]{Demailly2018AlmostKE}, supplies a positive
K\"ahler--Einstein metric on $X$.
Write its K\"ahler form as $\sigma$, with
$\rho(\sigma)=\lambda\sigma$ and $\lambda>0$. Constant rescaling
does not change the Ricci form, so $\omega=\lambda\sigma$ satisfies
$\rho(\omega)=\omega$. Since the Ricci form represents
$2\pi c_1(-K_X)$ in this convention, the asserted class follows.
\end{proof}

\subsection{Albanese fibres and the ordinary universal cover}

We now combine the classification with the analytic structure theorems
for compact K\"ahler manifolds with nef tangent bundle.

\begin{proof}[Proof of Corollary~\ref{cor:nef-tangent-uniformization}]
If $\dim_{\mathbb C}X=0$, connectedness makes $X$ a point and every
assertion follows. Assume henceforth that $\dim_{\mathbb C}X>0$.
Every finite \'{e}tale cover remains compact K\"ahler, and its tangent
bundle is analytically nef by Proposition~3.4 of
Demailly--Peternell--Schneider.
Their Proposition~3.9
\cite{DemaillyPeternellSchneider1994}, applied to every connected finite
\'{e}tale cover, makes its Albanese map surjective; hence its
irregularity is at most $\dim_{\mathbb C}X$. These irregularities are
nonnegative integers and include that of $X$, so their maximum
$q=\widetilde q(X)$ is attained.
Choose a connected finite \'{e}tale cover $X'\to X$ attaining it.
The Main Theorem and Theorem~3.14 of
Demailly--Peternell--Schneider, together with their Proposition~3.9,
give a smooth proper Albanese surjection
\[
 \alpha:X'\longrightarrow A=\operatorname{Alb}(X')
\]
with connected fibres, where $A$ is a complex torus of dimension $q$.
Every positive-dimensional fibre is a smooth projective Fano manifold
with nef tangent bundle and has dimension $d=\dim_{\mathbb C}X-q$.
Analytic nefness agrees with numerical nefness on these projective
fibres, so the algebraic classification applies.

If $d=6$, Theorem~\ref{thm:main} makes every fibre rational homogeneous.
For $1\leq d\leq5$, a projective fibre $F_a$ has
$\rho(F_a)\geq1>d-5$, and Kanemitsu's Theorem~0.2
\cite{Kanemitsu2016} gives the same conclusion. This is precisely the
established lower-dimensional range recalled in the introduction.
If $d=0$, each connected fibre is a point. This proves the first
part of the corollary.

Suppose first that $d>0$. Every rational homogeneous manifold $G/P$ satisfies
\[
 H^1(G/P,T_{G/P})=0
\]
and is locally rigid under complex deformations
\cite[Introduction and Section~3.7]{BienBrion1996}.
Applying this at each fibre of the proper
holomorphic submersion $\alpha$ shows that the fibre's biholomorphism
type is locally constant on $A$. Since $A$ is connected, all fibres
are biholomorphic to one fixed rational homogeneous manifold $F$.
The analytic Fischer--Grauert theorem
\cite{FischerGrauert1965}, in the form recalled in
\cite[Theorem~1.1]{Poczobut2025}, now makes $\alpha$ holomorphically
locally trivial. If $d=0$, the proper submersion with one-point fibres
is a biholomorphism, so this conclusion also holds with $F$ a point.

The universal cover of the complex torus $A$ is $\mathbb C^q$, which
is Stein and contractible even when $A$ is not projective.
Theorem~4.4 of Claudon--H\"oring--Koll\'ar
\cite{ClaudonHoringKollar2013} applies to the locally trivial proper
fibration $\alpha$ and gives
\[
 \widetilde {X'}\simeq\widetilde F\times\mathbb C^q.
\]
Here the tildes denote ordinary universal covers. A rational homogeneous
manifold $F$ is simply connected: write it using the simply connected
semisimple cover of its acting group; the corresponding parabolic
subgroup is connected, and the homotopy exact sequence of that
homogeneous fibration gives $\pi_1(F)=0$. This is also true for a point.
Thus $\widetilde F=F$. The connected finite \'{e}tale cover $X'\to X$
has the same ordinary universal cover as $X$, proving
$\widetilde X\simeq F\times\mathbb C^q$.

Since $F$ is projective, $F\times\mathbb A^q$ is quasi-projective and
its analytification is $F\times\mathbb C^q$. The projection of this
product to $\mathbb C^q$ is proper, surjective, and has connected compact
fibres. For every open set $U\subseteq\mathbb C^q$, a holomorphic
function on $F\times U$ is constant on each fibre and hence descends
holomorphically to $U$ by restriction to $\{f_0\}\times U$ for any
fixed $f_0\in F$. Thus the
projection has direct image of its structure sheaf equal to
$\mathcal O_{\mathbb C^q}$. Since the target is Stein, this projection
is the Remmert reduction, and the product is holomorphically convex.
The argument includes $q=0$, when the target is a point.
\end{proof}
